\section{A rerun-corrected upper bound}\label{sec:upperbound}

\looseness=-1 Appendix Figure~\ref{fig:ceilings} shows the two ceilings a perfect per-task choice could reach (per-cell values in Appendix Table~\ref{tab:t42}), and only one survives its own control. The \emph{success-rate ceiling} (any mode solves it) runs 7.1--51.9\% against a best single mode of 2.2--35.6\%; but it is a six-arm union quoted against a one-arm baseline. Matched at one added arm, the two purchases are: \textbf{a genuinely different representation buys +1.97 to +7.14pp; a second run of an arm already in hand buys 4.91--7.59pp} wherever a replicate exists. These are the same quantity at the same arm count in the same cell, so they are directly comparable, and repetition is not the cheaper of the two by accident: it is the larger. The \emph{cost ceiling at matched success} adds no arm and is untouched by that objection: keeping the best mode everywhere and sending only never-solved tasks to the cheapest arm leaves success unchanged by construction and cuts cost 9.5--30.6\% in 8 of 8 cells. The corresponding decomposition appears in Appendix Tables~\ref{tab:pb01}--\ref{tab:pb02}.

The same rerun correction applies to fusion as a representation purchase: Appendix Table~\ref{tab:t18} puts a distinct arm and a rerun side by side at the same arm count, and against the workload-matched single channel the fused mode's premium clears the rerun threshold of \S\ref{sec:noise} in 0 of 8 cells (Appendix Figure~\ref{fig:fusion} and Table~\ref{tab:t16}).

\looseness=-1 The number worth aiming at is therefore not the six-arm union. It is the gap between the fixed policies of \S\ref{sec:lowerbound} and the \emph{cost} ceiling: 9.5--30.6\% in 8 of 8 cells at unchanged success, reachable in principle with one bit per task rather than a mode identity. That target survives the rerun objection for the reason the success-rate ceiling does not: it adds no arm, being the same arms spent differently.
